% Figura: hmi-layout -- requiere % =====================================================================
%  Estilos compartidos por los diagramas TikZ de los capitulos.
%  Va en el PREAMBULO del documento principal:
%      % =====================================================================
%  Estilos compartidos por los diagramas TikZ de los capitulos.
%  Va en el PREAMBULO del documento principal:
%      % =====================================================================
%  Estilos compartidos por los diagramas TikZ de los capitulos.
%  Va en el PREAMBULO del documento principal:
%      \input{diagramas/estilos-diagramas}
% =====================================================================

\usepackage{tikz}
\usepackage{amsmath}
\usepackage{amssymb}
\usetikzlibrary{arrows.meta,positioning,shapes.geometric,shapes.misc,calc,fit,backgrounds,decorations.pathreplacing}

\definecolor{cbloque}{HTML}{E8EEF4}
\definecolor{cborde}{HTML}{4A6D8C}
\definecolor{cacento}{HTML}{D9E6D5}
\definecolor{cacentob}{HTML}{5A7A4A}
\definecolor{cgris}{HTML}{EDEDED}
\definecolor{cgrisb}{HTML}{888888}
\definecolor{calarma}{HTML}{F6DEDE}
\definecolor{calarmab}{HTML}{A03030}

\tikzset{
  bloque/.style={draw=cborde,fill=cbloque,rounded corners=2pt,align=center,
                 inner sep=5pt,font=\footnotesize,minimum height=9mm},
  bloqueg/.style={bloque,draw=cgrisb,fill=cgris},
  bloquea/.style={bloque,draw=cacentob,fill=cacento},
  bloquer/.style={bloque,draw=calarmab,fill=calarma},
  proceso/.style={bloque,minimum width=52mm},
  decision/.style={draw=cborde,fill=cbloque,align=center,font=\scriptsize,
                   diamond,aspect=2.2,inner sep=1pt},
  flecha/.style={-{Stealth[length=2mm]},draw=cborde!80!black,thick},
  flechag/.style={flecha,draw=cgrisb},
  et/.style={font=\scriptsize,inner sep=1.5pt,fill=white,align=center},
  etl/.style={et,midway},
  grupo/.style={draw=cgrisb,dashed,rounded corners=3pt,inner sep=4mm},
  tit/.style={font=\scriptsize\bfseries,text=cgrisb!60!black},
}

% --- utilidades para diagramas de secuencia ---
\tikzset{
  vida/.style={draw=cgrisb,dash pattern=on 2pt off 2pt},
  cab/.style={bloque,minimum width=26mm,font=\scriptsize\bfseries},
  nota/.style={draw=cacentob,fill=cacento,font=\scriptsize,align=center,
               rounded corners=1pt,inner sep=3pt},
}
\newcommand{\msg}[4]{\draw[flecha] (#1,#3) -- (#2,#3) node[et,midway,above=0pt]{#4};}
\newcommand{\msgr}[4]{\draw[flecha,dashed] (#1,#3) -- (#2,#3) node[et,midway,above=0pt]{#4};}
\newcommand{\auto}[3]{%
  \draw[flecha] (#1,#2) -- ++(0.55,0) -- ++(0,-0.42) -- (#1,{#2-0.42});
  \node[et,anchor=west] at ({#1+0.7},{#2-0.21}) {#3};}


% =====================================================================

\usepackage{tikz}
\usepackage{amsmath}
\usepackage{amssymb}
\usetikzlibrary{arrows.meta,positioning,shapes.geometric,shapes.misc,calc,fit,backgrounds,decorations.pathreplacing}

\definecolor{cbloque}{HTML}{E8EEF4}
\definecolor{cborde}{HTML}{4A6D8C}
\definecolor{cacento}{HTML}{D9E6D5}
\definecolor{cacentob}{HTML}{5A7A4A}
\definecolor{cgris}{HTML}{EDEDED}
\definecolor{cgrisb}{HTML}{888888}
\definecolor{calarma}{HTML}{F6DEDE}
\definecolor{calarmab}{HTML}{A03030}

\tikzset{
  bloque/.style={draw=cborde,fill=cbloque,rounded corners=2pt,align=center,
                 inner sep=5pt,font=\footnotesize,minimum height=9mm},
  bloqueg/.style={bloque,draw=cgrisb,fill=cgris},
  bloquea/.style={bloque,draw=cacentob,fill=cacento},
  bloquer/.style={bloque,draw=calarmab,fill=calarma},
  proceso/.style={bloque,minimum width=52mm},
  decision/.style={draw=cborde,fill=cbloque,align=center,font=\scriptsize,
                   diamond,aspect=2.2,inner sep=1pt},
  flecha/.style={-{Stealth[length=2mm]},draw=cborde!80!black,thick},
  flechag/.style={flecha,draw=cgrisb},
  et/.style={font=\scriptsize,inner sep=1.5pt,fill=white,align=center},
  etl/.style={et,midway},
  grupo/.style={draw=cgrisb,dashed,rounded corners=3pt,inner sep=4mm},
  tit/.style={font=\scriptsize\bfseries,text=cgrisb!60!black},
}

% --- utilidades para diagramas de secuencia ---
\tikzset{
  vida/.style={draw=cgrisb,dash pattern=on 2pt off 2pt},
  cab/.style={bloque,minimum width=26mm,font=\scriptsize\bfseries},
  nota/.style={draw=cacentob,fill=cacento,font=\scriptsize,align=center,
               rounded corners=1pt,inner sep=3pt},
}
\newcommand{\msg}[4]{\draw[flecha] (#1,#3) -- (#2,#3) node[et,midway,above=0pt]{#4};}
\newcommand{\msgr}[4]{\draw[flecha,dashed] (#1,#3) -- (#2,#3) node[et,midway,above=0pt]{#4};}
\newcommand{\auto}[3]{%
  \draw[flecha] (#1,#2) -- ++(0.55,0) -- ++(0,-0.42) -- (#1,{#2-0.42});
  \node[et,anchor=west] at ({#1+0.7},{#2-0.21}) {#3};}


% =====================================================================

\usepackage{tikz}
\usepackage{amsmath}
\usepackage{amssymb}
\usetikzlibrary{arrows.meta,positioning,shapes.geometric,shapes.misc,calc,fit,backgrounds,decorations.pathreplacing}

\definecolor{cbloque}{HTML}{E8EEF4}
\definecolor{cborde}{HTML}{4A6D8C}
\definecolor{cacento}{HTML}{D9E6D5}
\definecolor{cacentob}{HTML}{5A7A4A}
\definecolor{cgris}{HTML}{EDEDED}
\definecolor{cgrisb}{HTML}{888888}
\definecolor{calarma}{HTML}{F6DEDE}
\definecolor{calarmab}{HTML}{A03030}

\tikzset{
  bloque/.style={draw=cborde,fill=cbloque,rounded corners=2pt,align=center,
                 inner sep=5pt,font=\footnotesize,minimum height=9mm},
  bloqueg/.style={bloque,draw=cgrisb,fill=cgris},
  bloquea/.style={bloque,draw=cacentob,fill=cacento},
  bloquer/.style={bloque,draw=calarmab,fill=calarma},
  proceso/.style={bloque,minimum width=52mm},
  decision/.style={draw=cborde,fill=cbloque,align=center,font=\scriptsize,
                   diamond,aspect=2.2,inner sep=1pt},
  flecha/.style={-{Stealth[length=2mm]},draw=cborde!80!black,thick},
  flechag/.style={flecha,draw=cgrisb},
  et/.style={font=\scriptsize,inner sep=1.5pt,fill=white,align=center},
  etl/.style={et,midway},
  grupo/.style={draw=cgrisb,dashed,rounded corners=3pt,inner sep=4mm},
  tit/.style={font=\scriptsize\bfseries,text=cgrisb!60!black},
}

% --- utilidades para diagramas de secuencia ---
\tikzset{
  vida/.style={draw=cgrisb,dash pattern=on 2pt off 2pt},
  cab/.style={bloque,minimum width=26mm,font=\scriptsize\bfseries},
  nota/.style={draw=cacentob,fill=cacento,font=\scriptsize,align=center,
               rounded corners=1pt,inner sep=3pt},
}
\newcommand{\msg}[4]{\draw[flecha] (#1,#3) -- (#2,#3) node[et,midway,above=0pt]{#4};}
\newcommand{\msgr}[4]{\draw[flecha,dashed] (#1,#3) -- (#2,#3) node[et,midway,above=0pt]{#4};}
\newcommand{\auto}[3]{%
  \draw[flecha] (#1,#2) -- ++(0.55,0) -- ++(0,-0.42) -- (#1,{#2-0.42});
  \node[et,anchor=west] at ({#1+0.7},{#2-0.21}) {#3};}

 en el preambulo
\begin{tikzpicture}[x=1cm,y=1cm,
  reg/.style={draw=cborde,fill=cbloque,rounded corners=1pt},
  regg/.style={draw=cgrisb,fill=cgris,rounded corners=1pt},
  rega/.style={draw=calarmab,fill=calarma,rounded corners=1pt},
  titr/.style={font=\scriptsize\bfseries,align=center},
  sub/.style={font=\tiny,align=center,text=black!65}]

\draw[draw=cgrisb!70,fill=white,rounded corners=2pt] (-0.25,-0.25) rectangle (14.25,10.55);

% encabezado
\node[regg,minimum width=14cm,minimum height=0.9cm] (hdr) at (7,10.0) {};
\node[titr] at (7,10.0) {CONTROL DE LONGITUD DE ONDA HSRL};
\node[sub,anchor=west] at (0.1,10.0) {logo\\FCEFyN};
\node[sub,anchor=east] at (13.9,10.0) {logo\\SMN};

% barra de conexion
\node[reg,minimum width=14cm,minimum height=0.8cm] (con) at (7,8.95) {};
\node[sub,anchor=west] at (0.15,8.95)
 {PUERTO \fbox{COM3} $\circlearrowright$ \; BAUDIOS \fbox{115200} \; \fbox{Conectar}};
\node[sub,anchor=east] at (13.85,8.95)
 {$\bullet$ USB \; $\bullet$ SEEDER OK \; $\bullet$ ESCRITURAS OK \; MODO ÚNICO OK};

% banda de alarmas
\node[rega,minimum width=14cm,minimum height=0.65cm] (alm) at (7,8.0) {};
\node[sub,anchor=west,text=black] at (0.15,8.0)
 {$\triangle$ ALTA — TEMPERATURA FUERA DE LÍMITES (63,1400 °C, rango 62,50–63,10)};
\node[sub,anchor=east,text=black] at (13.85,8.0) {+1 más};

% panel izquierdo
\node[reg,minimum width=3.75cm,minimum height=7.5cm,anchor=north west] (pan) at (0.1,7.55) {};
\foreach \i/\t in {0/ESTADO DE CONTROL,1/TEMPERATURA HEATER,2/SINTONIZACIÓN AUTO,
                   3/VOLTAJE PIEZOELÉCTRICO,4/DETECTOR,5/SINTONIZACIÓN,6/CICLO}{
  \node[regg,minimum width=3.5cm,minimum height=0.95cm,anchor=north west]
        at (0.22,{7.42-\i*1.05}) {};
  \node[sub,anchor=north west] at (0.32,{7.36-\i*1.05}) {\t};
}

% grafico
\node[reg,minimum width=9.85cm,minimum height=3.9cm,anchor=north west] (plot) at (4.05,7.55) {};
\node[sub,anchor=north west] at (4.2,7.45) {AMPLITUD vs TEMPERATURA};
\draw[cgrisb] (4.85,4.0) -- (13.6,4.0);
\draw[cgrisb] (4.85,4.0) -- (4.85,6.9);
\foreach \x/\y in {5.2/4.9, 5.9/4.5, 6.6/4.3, 7.3/4.6, 8.0/5.4, 8.7/6.2, 9.4/6.6, 10.1/6.4}
  {\fill[cborde] (\x,\y) circle (1.3pt);}
\foreach \x/\y in {5.2/6.3, 5.9/6.6, 6.6/6.4, 7.3/5.6, 8.0/4.7, 8.7/4.3, 9.4/4.5, 10.1/5.2}
  {\fill[calarmab] (\x,\y) circle (1.3pt);}
\draw[cacentob,dashed] (7.65,4.05) -- (7.65,6.9);
\node[sub,anchor=south] at (7.65,6.9) {T. media};
\node[sub,anchor=south east] at (13.5,4.05) {Temperatura (°C) [PV]};
% ejes fijos: el vertical es todo el rango del ADC, el horizontal el rango del heater
\node[sub,anchor=east] at (4.78,4.0) {0};
\node[sub,anchor=east] at (4.78,6.9) {3,3 V};
\node[sub,anchor=north] at (4.85,3.95) {62,50};
\node[sub,anchor=north] at (13.6,3.95) {63,10};
\draw[cgrisb] (13.6,4.0) -- (13.6,4.12);

% controles
\node[reg,minimum width=9.85cm,minimum height=2.2cm,anchor=north west] (ctl) at (4.05,3.45) {};
\node[sub,anchor=north west] at (4.2,3.35) {CONTROLES};
\node[sub,anchor=north west,align=left] at (4.2,3.00)
 {\fbox{DETENER} \; \fbox{AVANZAR} \; \fbox{RETROCEDER} \; \fbox{SINTONIZAR} \; \fbox{SINTONIZAR AUTO}\\[1.5pt]
  Temp. SP \fbox{\;\;\;\;} \; Voltaje P. \fbox{\;\;\;\;} \; Promediado N \fbox{\;\;} \; Espera \fbox{\;\;} s\\[1.5pt]
  SEEDER \; \fbox{Consultar (!)} \fbox{Diagn. (d)} \fbox{Escribir SP (w)} \fbox{Desbloq. (u)} \fbox{Bloq. (l)}};

% monitor serie
\node[draw=cgrisb,fill=black!85,rounded corners=1pt,
      minimum width=9.85cm,minimum height=1.35cm,anchor=north west] (mon) at (4.05,1.15) {};
\node[sub,anchor=north west,text=white] at (4.2,1.05)
 {MONITOR SERIE\\[1pt]\ttfamily $\rightarrow$ g\\\ttfamily [ok] modo lock automatico activado\\\ttfamily [log] 27,g,62.8500,62.8471,\dots};

\node[nota,text width=82mm,anchor=north west] at (0,-0.9)
 {La banda de alarmas sólo se dibuja en condición anormal; en operación normal el
  cuerpo sube y ocupa su lugar. Los dos ejes del gráfico son fijos: el vertical cubre
  todo el rango del conversor ($0$--$3{,}3$\,V) y el horizontal el rango de seguridad
  del calefactor que declara el firmware.};

\end{tikzpicture}
